\documentclass[11pt,a4paper]{article}
\usepackage{fontspec,amsmath,amssymb,amsthm,microtype}
\setmainfont{lmroman10-regular.otf}[BoldFont=lmroman10-bold.otf,ItalicFont=lmroman10-italic.otf,BoldItalicFont=lmroman10-bolditalic.otf]
\setmonofont{lmmono10-regular.otf}[Scale=MatchLowercase]
\usepackage[margin=25mm]{geometry}
\usepackage[unicode,hidelinks]{hyperref}
\setlength{\emergencystretch}{3em}
\setlength{\parskip}{0.4em}
\setlength{\parindent}{0pt}
\allowdisplaybreaks[2]
\raggedbottom
\title{\texorpdfstring{Iteration of the scalar error estimate}{Iteration of the scalar error estimate}}
\author{LangevinLean mathematical exposition}
\date{}
\begin{document}\maketitle
\noindent Source module: \texttt{Unadjusted\allowbreak{}Langevin/\allowbreak{}Recurrence.\allowbreak{}lean}. The accompanying HTML pairs this mathematical exposition with the exact Lean statements and records the scope of the formalization.


Once coupling gives \(e_{k+1}\le qe_k+b\), the remaining calculation is deterministic. We retain the exact geometric sum before obtaining the simpler bias-floor bound. The upper envelope approaches \(b/(1-q)\); this does not assert that the actual error converges to that value.



The scalar sequence need not have an algorithmic interpretation for these lemmas. The proofs separate the exact geometric sum from its uniform bias-floor envelope and explicitly require the contraction factor in \([0,1)\). Later the coupling construction proves that the actual ULA error satisfies this recurrence.


\section{\texorpdfstring{Iteration gives a geometric sum}{Iteration gives a geometric sum}}


If \(q\ge0\) and \(e_{k+1}\le qe_k+b\) for every \(k\), then

\[e_n\le q^ne_0+b\sum_{j=0}^{n-1}q^j.\]

Induct on \(n\). At zero the sum is empty and both sides equal \(e_0\). Assume the bound at \(n\).



Multiply the induction inequality by \(q\ge0\), preserving its direction.



Write \(S_n=\sum_{j=0}^{n-1}q^j\). The geometric identity gives \(qS_n+1=S_n+q^n\); multiply this equality by \(b\).



Combine the one-step bound with the multiplied induction bound. Since \(S_{n+1}=S_n+q^n\), obtain \(e_{n+1}\le q^{n+1}e_0+bS_{n+1}\).



{\raggedright\noindent Lean declarations: \texttt{affine\_\allowbreak{}recurrence\_\allowbreak{}sum}.\par}

\section{\texorpdfstring{The exact finite-\allowbreak{}horizon envelope}{The exact finite-horizon envelope}}


If \(0\le q<1\) and \(e_{k+1}\le qe_k+b\), then

\[e_n\le q^ne_0+\frac{b(1-q^n)}{1-q}.\]

Apply the geometric-sum bound and evaluate the sum as \((q^n-1)/(q-1)\). Its denominator is nonzero because \(q<1\).



Change the signs of numerator and denominator simultaneously and put \(b\) in the numerator to obtain \(b(1-q^n)/(1-q)\).



{\raggedright\noindent Lean declarations: \texttt{affine\_\allowbreak{}recurrence\_\allowbreak{}exact}.\par}

\section{\texorpdfstring{A uniform bias-\allowbreak{}floor bound}{A uniform bias-floor bound}}


If \(0\le q<1\), \(b\ge0\), and \(e_{k+1}\le qe_k+b\), then

\[e_n\le q^ne_0+\frac b{1-q}.\]

Start from the exact finite-horizon bound. Its denominator \(1-q\) is positive.



Since \(q^n\ge0\) and \(b\ge0\), the numerator satisfies \(b(1-q^n)\le b\).



Divide by the positive denominator, add \(q^ne_0\), and compose with the exact bound.



{\raggedright\noindent Lean declarations: \texttt{affine\_\allowbreak{}recurrence\_\allowbreak{}bound}.\par}

\section{\texorpdfstring{Limit of the upper envelope}{Limit of the upper envelope}}


For \(0\le q<1\) and arbitrary real \(b,e_0\),

\[q^ne_0+\frac b{1-q}\longrightarrow\frac b{1-q}\quad(n\to\infty).\]

The geometric powers satisfy \(q^n\to0\).



Multiply this limit by \(e_0\) and add the constant \(b/(1-q)\).



{\raggedright\noindent Lean declarations: \texttt{bias\_\allowbreak{}envelope\_\allowbreak{}tendsto}.\par}

\section{\texorpdfstring{Every accuracy above the floor is eventually attained}{Every accuracy above the floor is eventually attained}}


Assume \(0\le q<1\), \(b\ge0\), and \(e_{k+1}\le qe_k+b\). For every \(\varepsilon>b/(1-q)\), the errors satisfy \(e_n<\varepsilon\) for every sufficiently large \(n\).



The upper envelope tends to a value strictly below \(\varepsilon\), so it is eventually below \(\varepsilon\).



At each such index, place \(e_n\) below the envelope by the recurrence bound and then below \(\varepsilon\).



{\raggedright\noindent Lean declarations: \texttt{eventually\_\allowbreak{}below\_\allowbreak{}bias}.\par}

\end{document}
